\chapter{Temporal Pattern Matching}
\label{ch:matcher}

The matcher is the engine's hot path. Every hunt the analyst runs,
every \code{run\_hunt} call from Claude, every CLI invocation ends
up here. The implementation ships three coordinated variants:

\begin{itemize}
    \item \textsc{Search-Temporal-Pattern} --- the entry point;
        dispatches work to rayon workers and collects results.
    \item \textsc{Dfs-Match-Iterative} --- the baseline iterative
        DFS executed per starting vertex.
    \item \textsc{Dfs-Match-Smart} --- an anomaly-bounded variant
        that maintains a top-$K$ min-heap and prunes branches
        whose running anomaly estimate makes them inadmissible.
\end{itemize}

A fourth path, \textsc{Search-Temporal-Pattern-Simd}, delegates to
\code{libgraphmatch} or the pure-Rust \code{simd\_rust} module
(Chapter~\ref{ch:simd}) and is gated on the \code{simd} feature
plus a runtime CPU probe.

\section{Shared state}

\begin{lstlisting}[style=rust]
// Scratch allocated once per worker and reused across starts.
struct DfsFrame {
    sid:          StrId,
    step_idx:     usize,
    edge_hi:      usize,   // upper bound of outgoing edges
    edge_cursor:  usize,   // next edge index to try
}

struct SmartDfsFrame {
    sid:              StrId,
    step_idx:         usize,
    edge_hi:          usize,
    edge_cursor:      usize,
    path_anomaly_sum: f64,   // sum of node_anomaly_estimate along path
}
\end{lstlisting}

Stack frames are plain structs stored in a \code{Vec<Frame>} that is
reused across starting vertices. This is the single biggest
allocation savings in the matcher --- the old recursive
implementation allocated per-frame on the system stack and could
not trivially restart.

\begin{invariant}[$k$-simplicity]
The running \code{visit\_count: HashMap<StrId, usize>} maintained
inside every DFS walk satisfies \code{visit\_count[v] <= k} for
every $v$ on the current path, where $k$ is
\code{hypothesis.k\_simplicity}. Violating this would admit paths
that revisit a vertex more than $k$ times, contradicting the
hypothesis' semantic. The algorithm enforces the invariant at
push-time: if incrementing \code{visit\_count[dest]} would exceed
$k$, the edge is skipped.
\end{invariant}

\section{Entry point: \textsc{Search-Temporal-Pattern}}

The outer algorithm partitions work by starting vertex. For each
entity whose type matches $H$'s first step, a rayon-managed worker
thread runs \textsc{Dfs-Match-Iterative} on that starting vertex
with its own scratch stack and result buffer. When every worker has
finished, results are merged, scored, deduplicated and paginated
(Chapter~\ref{ch:dedup-pagination}).

\begin{algorithm}
\caption{\textsc{Search-Temporal-Pattern}}\label{alg:search-entry}
\KwIn{graph $G$, hypothesis $H$, optional window $(t_0, t_1)$,
      optional cap $\mathrm{max}$}
\KwOut{vector of path vectors}
$\mathrm{starts} \gets $ \code{G.type\_index[H.steps[0].origin\_type]}\;
\If{$\mathrm{starts} = \emptyset$}{\Return $\emptyset$\;}
$\mathrm{cap} \gets $ \code{AtomicUsize(0)}\;
$\mathrm{results} \gets $ parallel-map over $\mathrm{starts}$:
    \ForEach{$s \in \mathrm{starts}$ \KwIn{parallel}}{
        $\mathrm{local} \gets $ \textsc{Dfs-Match-Iterative}($G, H, s, t_0, t_1, \mathrm{cap}, \mathrm{max}$) \;
        \Return $\mathrm{local}$\;
    }
\Return concat of $\mathrm{results}$\;
\end{algorithm}

The \code{AtomicUsize} cap is read by every worker on every
candidate push so that once the aggregate result count reaches
$\mathrm{max}$, remaining workers stop as early as possible. This
is a classical early-exit pattern: the count may briefly exceed
$\mathrm{max}$ because workers commit then check, but the overshoot
is bounded by the number of rayon workers (typically 8--16).

\section{The iterative DFS}

\begin{algorithm}
\caption{\textsc{Dfs-Match-Iterative}}\label{alg:dfs-iter}
\KwIn{graph $G$, hypothesis $H$, start $s_0$, window $(t_0, t_1)$,
      shared cap counter $\mathrm{cap}$, cap $\mathrm{max}$}
\KwOut{list of matching paths rooted at $s_0$}
$\mathrm{results} \gets \emptyset$;
$\mathrm{stack} \gets \emptyset$;
$\mathrm{path} \gets [s_0]$;
$\mathrm{visit} \gets \{s_0 : 1\}$\;
\textsc{Push-Edges-For-Step}($\mathrm{stack}$, $G$, $s_0$, step $0$, $t_0, t_1$)\;
\While{$\mathrm{stack} \neq \emptyset$}{
    $f \gets $ top of $\mathrm{stack}$ \;
    \If{$f.\text{edge\_cursor} \geq f.\text{edge\_hi}$}{
        pop $f$;
        pop last vertex from $\mathrm{path}$;
        decrement $\mathrm{visit}[\cdot]$;
        \Continue\;
    }
    $c \gets $ \code{G.edge\_store[f.sid][f.edge\_cursor]}\;
    $f.\text{edge\_cursor} \gets f.\text{edge\_cursor} + 1$\;
    \tcp{filter: relation type}
    \If{$c.\text{rel\_type\_tag} \neq $ step[$f.\text{step\_idx}$].rel}{
        \Continue\;
    }
    \tcp{filter: destination entity type via SoA tag}
    $u \gets c.\text{dest\_sid}$\;
    \If{$\code{G.entity\_type\_tags}[u.0] \neq $ step[$f.\text{step\_idx}$].dst}{
        \Continue\;
    }
    \tcp{filter: time window}
    \If{$c.\text{timestamp} \notin [t_0, t_1]$}{\Continue\;}
    \tcp{filter: edge predicates (on metadata)}
    \If{\textsc{Evaluate-Predicates}(step.edge\_predicates, $c$, $G$) = false}{
        \Continue\;
    }
    \tcp{cycle prevention via $k$-simplicity}
    \If{$\mathrm{visit}[u] + 1 > H.k$}{\Continue\;}
    $\mathrm{visit}[u] \gets \mathrm{visit}[u] + 1$;
    append $u$ to $\mathrm{path}$\;
    \If{$f.\text{step\_idx} + 1 = $ len($H.\text{steps}$)}{
        record $\mathrm{path}$ in $\mathrm{results}$\;
        increment $\mathrm{cap}$ atomically; \If{$\mathrm{cap} \geq \mathrm{max}$}{break loop\;}
        $\mathrm{visit}[u] \gets \mathrm{visit}[u] - 1$;
        pop last vertex;
        \Continue\;
    }
    push new frame: $(u, f.\text{step\_idx}+1, \ldots)$ onto $\mathrm{stack}$\;
    \textsc{Push-Edges-For-Step}($\mathrm{stack}$, $G$, $u$, $f.\text{step\_idx}+1$, $t_0, t_1$)\;
}
\Return $\mathrm{results}$\;
\end{algorithm}

\textsc{Push-Edges-For-Step} is a one-liner that binary-searches
\code{edge\_store[v]} for the first edge with \code{timestamp}
$\geq t_0$ and sets \code{edge\_hi} to the first edge with
\code{timestamp} $> t_1$. The sort by timestamp performed at
finalize (Chapter~\ref{ch:graph-hunter}) is what makes this binary
search correct.

\begin{theorem}[Matcher correctness]
Upon return, \textsc{Dfs-Match-Iterative} produces exactly the set
of temporal paths of length $\ell$ rooted at $s_0$ that satisfy the
hypothesis $H$ within $[t_0, t_1]$ and respect $k$-simplicity.
\end{theorem}

\begin{proof}[Proof sketch]
By induction on the hypothesis depth. The base case ($\ell = 1$)
is the first push into \code{stack}; the check is step-by-step
equivalent to the hypothesis constraints on the first edge. The
inductive step observes that each admitted edge extends the path by
exactly one vertex, updates \code{visit} consistently, and that
the cycle check rejects exactly those extensions that would exceed
$k$-simplicity. The stack unwinds restoring \code{visit} and the
path, so every branch is explored independently; no path is ever
omitted because of missed backtracking. \qed
\end{proof}

\begin{complexity}
The worst case is a graph where every edge matches every step.
Then from a start $s_0$ the DFS visits up to
$\bar{b}^\ell$ paths (where $\bar{b}$ is the average out-degree of
qualifying edges per step and $\ell$ is the step count), giving
$\Ord{\bar{b}^\ell \cdot \ell}$ time per start. Across all $n_s$
candidate starts, the total is
\[
    T_{\text{match}}(G, H) = \Ord{n_s \cdot \bar{b}^\ell \cdot \ell}.
\]
Space per worker is $\Ord{\ell}$ for the stack and path plus
$\Ord{\ell \cdot k}$ for the \code{visit} map (each entry is one
vertex, and the path holds at most $\ell$ distinct vertices each
visited at most $k$ times). With $W$ workers total matcher space is
$\Ord{W \cdot \ell \cdot k}$. In practice $\bar{b}$ is tiny because
the type-tag fast-path eliminates most outgoing edges on the first
byte comparison; empirical results from \code{benches/hunt\_latency.rs}
are in App.~\ref{app:bench}.
\end{complexity}

\section{The smart variant}

\textsc{Dfs-Match-Smart} is used when the caller requests the
top-$K$ most anomalous paths instead of ``all paths.'' It differs
from the iterative variant only in the candidate ordering and the
result accumulator:

\begin{itemize}
    \item \textbf{Accumulator}: a min-heap of size at most $K$
        ordered by composite anomaly score. When full, new paths
        are accepted iff they score higher than the heap's minimum.
    \item \textbf{Running bound}: each frame tracks
        \code{path\_anomaly\_sum}, the sum of
        \code{node\_anomaly\_estimate} (Chapter~\ref{ch:anomaly}) along
        the path so far. Extending the path cannot decrease this
        sum; if the sum already falls below the heap's minimum,
        the branch is pruned.
\end{itemize}

\begin{algorithm}
\caption{\textsc{Merge-Into-TopK}}\label{alg:merge-topk}
\KwIn{local heap $L$, global heap $G$, capacity $K$}
\ForEach{$p \in L$}{
    \If{$|G| < K$}{$G$.push($p$)}
    \ElseIf{$\text{score}(p) > \min(G)$}{
        $G$.pop\_min(); $G$.push($p$);
    }
}
\end{algorithm}

\begin{complexity}
$\Ord{(|L| + |G|) \log K}$. When the matcher joins rayon workers,
each local heap of size $\leq K$ merges into a single global heap
at total cost $\Ord{W K \log K}$, negligible compared to the DFS
work.
\end{complexity}

\section{Parallelism}

The matcher is embarrassingly parallel in the starting vertex.
Because each worker's DFS reads the graph exclusively (all writers
are blocked on an \code{RwLock::write} at the session level), there
is no inter-worker synchronisation beyond the shared \code{cap}
counter and the final reduction into a single result vector.

\begin{theorem}[Parallel speedup]
Given $W$ rayon workers and $n_s$ starting vertices with
approximately equal DFS work per start, the matcher achieves a
speedup factor of $\min(W, n_s)$ up to the constant cost of the
final reduction.
\end{theorem}

The proof is Brent's theorem applied to the trivial partition: each
start's DFS has no inter-start dependencies, and the final
reduction is $\Ord{\text{total results}}$ which is typically orders
of magnitude smaller than the DFS itself.

\begin{inthecode}
\filepath{graph\_hunter\_core/src/graph.rs:461} ---
\code{search\_temporal\_pattern}.\\
\filepath{graph\_hunter\_core/src/graph.rs:1121} ---
\code{dfs\_match\_iterative}.\\
\filepath{graph\_hunter\_core/src/graph.rs:782} ---
\code{search\_temporal\_pattern\_smart}.\\
\filepath{graph\_hunter\_core/src/graph.rs:932} ---
\code{dfs\_match\_smart}.\\
\filepath{graph\_hunter\_core/benches/hunt\_latency.rs} ---
performance benchmarks.
\end{inthecode}
